\chapter{Augmented Reality}

Augmented Reality (AR) is a technology that superimposes computer-generated content onto the real world, enhancing the user's perception and interaction with their environment. There are two different types of environments:
\begin{itemize}
    \item Real Environment: The physical world that we live in, which can be seen and interacted with directly.
    \item Virtual Environment: A computer-generated environment that may (or may not) resemble the real world, and is typically experienced through a device such as a headset or smartphone.
\end{itemize}

When the virtual and real environments are combined, we get what is known as Mixed Reality (MR). There are two different types of MR:
\begin{itemize}
    \item Augmented Reality (AR): The virtual content is overlaid onto the real world, allowing users to see and interact with both the real and virtual elements simultaneously.
    \item Augmented Virtuality (AV): The real world is integrated into a virtual environment, allowing users to interact with real-world objects within a virtual space.
\end{itemize}

\section{Devices for Augmented Reality}

In order to experience Augmented Reality, users typically require specific devices that can capture the real world and display the augmented content. Some common devices for AR include Smartphones, Tablets, AR Glasses, and Head-Mounted Displays (HMDs). Each of these devices has its own advantages and limitations in terms of portability, field of view, and interaction capabilities. Some aspects to consider when choosing an AR device include:
\begin{description}
    \item[Video vs Optical See-Through]~
    \begin{itemize}
        \item  \ul{Video see-through (VST)}: devices use cameras to capture the real world, and then computer-generated content is superimposed onto the video feed.
        
        It can be implemented using a smartphone or tablet, where the device's camera captures the real world and the screen displays the augmented content. It also allows a full peripheral view of the real world, which may be needed in dangerous situations.
        \item \ul{Optical see-through (OST)}: devices have two layers: a transparent layer that allows users to see the real world directly, and a reflective layer that projects virtual content onto the transparent layer.
        
        OST devices typically require specialized hardware such as AR glasses or HMDs, which can be more expensive and less portable than VST devices. It can also provide a more immersive experience, but may suffer from latency and lower image quality compared to VST. Lastly, users may suffer from motion sickness or eye strain when using OST devices for extended periods of time.
    \end{itemize}

    \item[HMD vs Mobile vs Projector]~
    
    There are different specific devices for AR, each with its own advantages and disadvantages. Each one has a spatial relationship to the user and the real environment, which needs parameters to be known. There are two types of parameters:
    \begin{itemize}
        \item \ul{C}: \textbf{C}onstant and \textbf{C}alibrated parameters that remain stable during usage.
        \item \ul{T}: \textbf{T}racked parameters that change constantly and therefore need tracking.
    \end{itemize}
    
    Depending on the device, the parameters that need to be tracked may differ.
    \begin{itemize}
        \item \ul{HMDs used for OST}: The distance to the real world needs to be tracked (T), while the distance to the eyes is constant (C). Depending if it uses eye-tracking or not, that mey be tracked (T) or constant (C).
        \item \ul{HMDs used for VST}: The distance to the real world needs to be tracked (T), and the distance between the camera and the display is constant (C), as they are integrated. The distance to the eyes is either constant (C) or tracked (T), depending on whether the device uses eye-tracking or not.
        \item \ul{Mobile devices}: The distance to the real world needs to be tracked (T), while the distance from the camera to the device is constant (C), as they are integrated. The distance to the eyes usually does not matter.
        \item \ul{Projectors}: The eyes do not matter, as the content is projected onto a surface. The distance to the real world is either constant (C) or tracked (T), depending on whether the scene is static or dynamic.
    \end{itemize}

    \item[Occularity in HMDs]~
    \begin{itemize}
        \item \ul{Monocular}: The virtual content is displayed to only one eye, which can lead to a less immersive experience but is cheaper and less complex to implement.
        \item \ul{Bi-ocular}: The virtual content is displayed to both eyes, but the same image is shown to each eye. This can provide a more immersive experience than monocular devices, but may still lack depth perception.
        \item \ul{Binocular}: The virtual content is displayed to both eyes, with each eye receiving a slightly different image to create a sense of depth and immersion. This can provide the most immersive experience, but is also the most complex and expensive to implement.
    \end{itemize}

    \item[Distance to the display]~
    \begin{itemize}
        \item \ul{Head Space}: The display is located close to the user's head, such as in HMDs.
        \item \ul{Body Space}: The display is located at a distance from the user's body, such as in mobile or handheld devices.
        \item \ul{World Space}: The display is located in the environment, such as stationary displays or projectors.
    \end{itemize}

    \item[Field of View (FOV)]~
    
    The FOV is the extent of the observable world that can be seen at any given moment. A wider FOV can provide a more immersive experience, but may also require more processing power and higher-quality hardware to maintain image quality and reduce latency. Humans real FOV is usually larger than most AR devices can provide, which can lead to a less immersive experience and may cause users to feel disconnected from the augmented content.
\end{description}

\section{Making sense of the real environment}

In order to effectively augment the real environment, AR systems need to be able to understand and interpret the physical world. This involves several key components, as registration, tracking, and calibration.

\subsection{Registration}

Registration is the process of aligning the virtual content with the real world in a way that appears seamless and natural to the user. In order to do so, the real-world coordinate systems need to be mapped to the virtual coordinate systems. Several transformations are needed to achieve this mapping, including:
\begin{itemize}
    \item \ul{Model Transformation}: This transformation maps the virtual content from its own coordinate system to a common world coordinate system.
    
    If the real objects move and the virtual content needs to be registered to them, tracking is needed to update the model transformation in real time.
    \item \ul{View Transformation}: This transformation maps the world coordinate system to the user's viewpoint, taking into account the user's position and orientation in the real world.
    
    If the camera is moving, tracking is needed to update the view transformation in real time.
    \item \ul{Projection Transformation}: This transformation maps the 3D virtual content onto the 2D display of the AR device, taking into account the device's field of view and other display parameters.
    
    Calibration is needed to determine the parameters of the projection transformation.
\end{itemize}

Depending on where the objects are registered, there are different types of registration:
\begin{itemize}
    \item \ul{Object Registration}: The virtual content is registered to specific objects in the real world, such as a table or a building. If the objects are moving, tracking is needed to update the registration in real time.
    \item \ul{View Registration}: The virtual content is registered to the user's viewpoint, so that it appears to be anchored in the real world regardless of the user's position or orientation. If the camera is moving, tracking is needed to update the registration in real time.
    \item \ul{World Registration}: The virtual content is registered to a fixed location in the real world, such as a specific point on the ground or a landmark. Regardless of whether the user or the objects are moving, no tracking is needed to maintain the registration, as it is anchored to a fixed location in the real world.
\end{itemize}

In addition, there are two main types of registration techniques, depending on how well objects are registered:
\begin{itemize}
    \item \ul{Geometric Registration}: it focuses on just placing the virtual content in the correct position and orientation in the real world.
    \item \ul{Photometric Registration}: it focuses on matching the lighting and shading of the virtual content to the real world, in order to create a more seamless and realistic integration.
    
    It is much more computationally expensive than geometric registration, as it requires real-time analysis of the lighting conditions in the real world and the ability to adjust the virtual content accordingly.
\end{itemize}

Lastly, \emph{occlusion} should be considered in registration, as it can affect the realism of the augmented content. If virtual objects are not properly occluded by real objects, it can break the illusion of integration and make the augmented content appear less believable. Proper occlusion handling requires accurate depth sensing and tracking of both the real and virtual objects in the environment.

\subsection{Tracking}

Tracking is the process of continuously monitoring and measuring the real-world environment and the user's position and orientation. There are different types of tracking systems, including:
\begin{itemize}
    \item \ul{Stationary Tracking}: The tracking system is fixed in a specific location and tracks the user's position and orientation relative to that location.
    \item \ul{Mobile Tracking}: The tracking system is integrated into the AR device and tracks the user's position and orientation relative to the real world.
    \item \ul{Optical Tracking}: The tracking system uses cameras and computer vision algorithms to track the user's position and orientation based on visual markers or features in the real world.
\end{itemize}

There are two main approaches to tracking in AR:
\begin{itemize}
    \item \ul{Model-Based Tracking}: This approach uses a pre-defined model of the real world, such as a 3D map or a set of known landmarks, to track the user's position and orientation. The tracking system compares the real-world environment to the model and updates the user's position and orientation accordingly.
    
    \item \ul{Model-Free Tracking}: This approach does not rely on a pre-defined model of the real world, but instead uses real-time sensor data to track the user's position and orientation.
\end{itemize}

Different specific tracking techniques include:
\begin{itemize}
    \item \ul{Marker-Based Tracking}: This technique uses visual markers, such as QR codes or fiducial markers, placed in the real world.
    
    \item \ul{Natural Feature Tracking}: This technique uses computer vision algorithms to track natural features in the real world, such as edges, corners, or textures.
    
    An specific example of this technique is Simultaneous Localization and Mapping (SLAM), which combines:
    \begin{itemize}
        \item Visual natural attribute tracking
        \item Depth sensing using specialized sensors.
    \end{itemize}
\end{itemize}

One specific object that usually needs to be correctly tracked is the user's hands and fingers, as they are often used for interaction with the virtual content.

\subsection{Calibration}

Calibration is the process of determining the parameters of the AR system, such as the position and orientation of the cameras, the field of view, and the display characteristics. Proper calibration is essential for accurate registration and tracking, as it ensures that the virtual content is correctly aligned with the real world and that the user's movements are accurately tracked.


\section{Situated AR}

An important aspect of AR is the concept of \emph{situated AR}, opposite to conventional AR.
\begin{itemize}
    \item \ul{Conventional Visualization}: The virtual content is displayed without integration with the real world. This includes traditional displays, or even AR content that is simply overlaid on the real world.
    
    This includes context registered relatively to markers, as they are added to the real world but do not have a specific meaning in it.
    \item \ul{Situated Visualization}: The virtual content is integrated into the real world in a way that is meaningful and relevant to the user's context.
\end{itemize}

Some situated AR techniques include:
\begin{itemize}
    \item \ul{Annotations and Labels}: Virtual content is used to provide additional information about real-world objects or locations, such as labels, descriptions, or instructions.
    
    They should not cover real-world objects nor overlap with other annotations.

    \item \ul{X-Ray Visualization}:
    
    This technique allows users to see through real-world objects, providing a view of hidden or obscured content. Transparency should still be used to keep up occlusion and depth impression.

    Some specific examples of x-ray visualization include showing the internal components of a machine, or allowing users to see through walls to detect wires or pipes.
    \item \ul{Spatial Manipulation}:
    
    This technique allows users to manipulate virtual content in the real world, such as moving, resizing, or rotating virtual objects. It can be used for tasks such as interior design, where users can visualize how furniture would look in a room before making a purchase.
    \item \ul{Information Filtering}:
    
    This technique allows users to filter or hide certain real-world objects or information, in order to focus on specific aspects of the environment. There are two main types of information filtering:
    \begin{itemize}
        \item \ul{Knowledge-Based Filtering}: This technique uses the user's knowledge and preferences to filter information.
        
        For instance, combined with X-Ray Visualization, it can be used to show only specific types of information, such as showing only the water pipes in a building while hiding the electrical wires.

        
        \item \ul{Spatial Filtering}: Use geometric and geographic information to place information. For instance, it can be used to show information about nearby points of interest while hiding information about distant locations.
    \end{itemize}
\end{itemize}


\section{AR Interaction}

Interaction in AR can be achieved through various input methods and output modalities.
\subsection{Output Modelities}

There are three main output modalities in AR:
\begin{itemize}
    \item \ul{Head-Referenced Display}: The virtual content is displayed in a way that is anchored to the user's FoV, such as in smart glasses.
    \item \ul{Hand-Referenced Display}: The virtual content is displayed in the user's palm of hand(s), allowing for direct interaction with the augmented content.
    \item \ul{Agile Display}: The virtual content is displayed everywhere and with changing perspectives of users, such as in mobile displays (carried around or worn) or stationary displays (adaptive projectors).
\end{itemize}

\subsubsection{Avatars}

Avatars are virtual representations of users in the AR environment, allowing for social interaction and communication with other users. Some different aspects that should be taken into account when designing avatars include:
\begin{itemize}
    \item Full or half body representation.
    \item Full featured or simplified abstract representation.
    
    Even though a full featured avatar can provide a more immersive experience and are usually more trusted, they can also be more expensive and complex to implement, and may require more processing power and higher-quality hardware to maintain image quality and reduce latency.
    \item Life size or scaled representation.
    
    Miniature avatars are usually preferred over life-size ones, and a potential reason for this is that they can fit in the user's FoV more easily.
\end{itemize}

\subsection{Input Modalities}

There are several input modalities for AR interaction, including:
\begin{itemize}
    \item \ul{Rigid Object Tracking}: This technique tracks the position and orientation of your head (HMD) or different devices (mobile, projectors) to allow for interaction with the virtual content.
    \item \ul{Gesture Recognition}: Using hand tracking, the system can recognize specific hand gestures to trigger interactions with the virtual content. This gesture recognition should be as intuitive and natural as possible.
    \item \ul{Pointing}: A laser beam or a ray can be emitted from the user's hand or device to point at the AR scene, allowing for selection of objects and non-verbal communication.
    \item \ul{Drawing}: Similar to pointing, but instead of just selecting objects, the user can draw in the air to create virtual content or annotate the real world.
    \item \ul{Touch}: As in real interfaces, virtual buttons or controls can be placed in the AR scene, allowing users to interact with them through touch.
\end{itemize}

\subsubsection{The Midas Touch Problem}

An important issue in AR interaction is selecting objects in the AR scene.
\begin{itemize}
    \item \ul{Selection via Gaze}
    
    A first approach to selection was using the user's gaze, where the system tracked the user's eye movements to determine where they were looking in the AR scene.

    However, this lead to the \emph{Midas\footnote{King Midas was a figure in Greek mythology who was granted the ability to turn everything he touched into gold. Even though this power was initially seen as a blessing, it quickly became a curse, as Midas could not eat, drink, or interact with anything without turning it into gold.}
    Touch Problem}, where every object that the user looked at would be selected, making it difficult to interact with the AR content without accidentally selecting things.

    \item \ul{Selection via Pointing or Gestures}
    
    In order to solve the Midas Touch Problem, selection can be done through pointing or gestures, where the user explicitly indicates which object they want to interact with.
\end{itemize}